\subsection{Dataset: Multi-Omics Conflict Benchmark}
The protocol was evaluated on a synthetic but high-fidelity dataset comprising:
\begin{itemize}
    \item \textbf{Proteomics}: 5,200 protein expression samples.
    \item \textbf{Genomics}: 12,400 SNP variants.
    \item \textbf{Conflicts}: 850 logical inconsistencies (e.g., conflicting P-values across datasets) designed to trigger hallucinations.
\end{itemize}

\subsection{Baseline Methods}
We compared the protocol against: (1) Genetic Programming (DEAP); (2) PPO-based RL optimization; (3) Constitutional AI with static principles; and (4) Static multi-agent frameworks (LangChain).

\subsection{Hardware and Digital Twin Environment}
The experimentation was conducted on an \textbf{Intel Arc A770 (XPU)} cluster (4 units) utilizing \textbf{oneAPI} acceleration, 128GB RAM, and Docker-based sandbox isolation.

\subsection{Computational Complexity Analysis}
\textbf{Algorithm 1 (Dialectical Protocol):} The total complexity per generation is $O(n \cdot d + k^2 + k \cdot m)$, where $n$ is lines of code, $d$ is AST depth, $k$ is mutation count, and $m$ is LLM latency. For a 500-line codebase, the average turnaround time is $\sim$45 seconds.

\textbf{Algorithm 2 (Genetic Proofreading):} Complexity is $O(n + g \cdot m)$, where $g$ is the number of gaps detected. This ensures linear scaling with code size during the enzymatic repair phase.
